% EM: Need a new name:
% Gradient-Conditioned Model Adapters

% Key ideas:
% - Editing a *pre-trained* model
% - Not architecture-specific (to e.g. transformers)
% - Scales well
% - Uses factorized form of gradient

%%% Paper titles?
% Editing pre-trained models with lossless gradient conditioning
% Editing large pre-trained models with lossless gradient conditioning
% Editing knowledge in large pre-trained models with lossless gradient conditioning
% Adapting knowledge in large pre-trained models with lossless gradient conditioning
% Modifying knowledge in large pre-trained models with lossless gradient conditioning
% Quickly Editing Large Pre-Trained Models
% A Fast Method for Editing Large Pre-Trained Models
% A Fast Method for Adapting Large-Scale Pre-Trained Models
% A Fast Method for Adapting Billion-Parameter Pre-Trained Models
% Model Editors for Adapting Large Pre-Trained Models
% Quick Fix: Model Editing at Scale 
% Quick Fix: Post-Hoc Model Editing at Scale 
% Quick Fix: Editing Pre-Trained Models at Scale 
% Editing Pre-Trained Models at Scale 


%%% Method names
% Lossless Gradient Conditioning
% Learnable Gradient Manipulation
% Knowledge Editing with Gradient Manipulation
% Knowledge Adaptation by Gradient Editing
% knowledge adapters
% scalable knowledge editors 
% scalable network editors (SNEs)
% scalable weight editors 
% lightweight model editor
%%% model editors
%% QuickFix

\DeclareMathOperator*{\argmax}{\arg\!\max}

\newcommand{\fullname}{Model Editor Networks with Gradient Decomposition}
\newcommand{\name}{MEND}
\newcommand{\baseparams}{\theta}
\newcommand{\editedparams}{\theta_e}
\newcommand{\xe}{x_\text{e}}
\newcommand{\xloc}{x_\text{loc}}
\newcommand{\yloc}{y_\text{loc}}
\newcommand{\ye}{y_\text{e}}
\newcommand{\losse}{L_\text{e}}
\newcommand{\lossl}{L_\text{loc}}
\newcommand{\losstotal}{L_\text{MEND}}
\newcommand{\dtrain}{D^{tr}_{edit}}
\newcommand{\dtest}{D^{te}_{edit}}
\newcommand{\dbase}{D_{base}}
\newcommand{\basemodel}{f_\baseparams}
\newcommand{\pseudograd}{edit update}
\newcommand{\neighborhood}{N(\xe,\ye)}
% \newcommand{\rebuttal}[1]{\textcolor{blue}{#1}}
\newcommand{\rebuttal}[1]{#1}
% \newcommand{\cameraready}[1]{\textcolor{orange}{#1}}
\newcommand{\cameraready}[1]{#1}

\newcommand{\mrow}[2]{\multirow{#1}{*}{\shortstack{#2}}}

\newcommand{\textlt}[1]{\textless\,#1\phantom{\textless\,}}
\newcommand{\textgt}[1]{\textgreater\,#1\phantom{\textgreater\,}}

% https://tex.stackexchange.com/questions/42619/x-mark-to-match-checkmark
\newcommand{\cmark}{\ding{51}}
\newcommand{\xmark}{\ding{55}}

\newcommand{\greenyes}{\textcolor{Green}{\ding{51}}}
\newcommand{\yellowmaybe}{\textcolor{YellowOrange}{\textbf{?}}}
\newcommand{\redno}{\textcolor{Red}{\ding{55}}}

% https://tex.stackexchange.com/a/385527/250175
\newcommand*{\fakemonochar}[1]{\makebox[.6em][c]{#1}}
\makeatletter
\newcommand*{\fakemonotext}[1]{%
  \@fakemonotext#1\\%
}
\newcommand*{\@fakemonotext}[1]{%
  \ifx #1\\\else\fakemonochar{#1}\expandafter\@fakemonotext\fi
}

% chris comments
\newif\ifcomments
\commentstrue
\ifcomments
\usepackage[normalem]{ulem}
\addtolength{\marginparwidth}{1in}
\definecolor{CMpurple}{rgb}{0.6,0.18,0.64}
\newcommand\cm[1]{\textcolor{CMpurple}{\textsf{\scriptsize[\textbf{CM\@:} #1]}}}
\newcommand\cmi[1]{\textcolor{CMpurple}{#1}}
\newcommand\cmm[1]{\marginpar{\raggedright\tiny\textcolor{CMpurple}{\textsf{{\bfseries CM\@:} #1}}}}
\newcommand\cms{\bgroup\markoverwith{\textcolor{CMpurple}{\rule[.4ex]{2pt}{0.8pt}}}\ULon}
\else
\newcommand\cm[1]{}
\newcommand\cmi[1]{\ignorespaces}
\newcommand\cmm[1]{}
\newcommand\cms[1]{#1}
\fi

\newcommand\antoine[1]{{\color{red}\{\textit{#1}\}$_{antoine}$}}
